\chapter{Manipulação de Dados}

\section{\texttt{generate} e \texttt{mutate}}

Criam ou substituem colunas. As duas formas são equivalentes em resultado; diferem
na semântica de mutação:

\begin{center}
\begin{tabular}{lll}
\toprule
\textbf{Forma} & \textbf{Semântica} & \textbf{Exemplo} \\
\midrule
\texttt{generate df col = expr}     & Imperativa (muta \texttt{df}) & estilo Stata \\
\texttt{let df2 = generate(df, col = expr)} & Funcional (df inalterado) & estilo funcional \\
\texttt{df |> generate(col = expr)} & Pipe assign-back (muta \texttt{df}) & encadeamento \\
\bottomrule
\end{tabular}
\end{center}

\begin{lstlisting}[caption={generate: forma imperativa}]
load "dados.csv" as df

generate df lrenda = log(renda)
generate df renda2 = renda^2
generate df dummy = if(renda > 1000, 1, 0)
\end{lstlisting}

\begin{lstlisting}[caption={mutate: forma funcional}]
let df2 = mutate(df, lrenda = log(renda), renda2 = renda^2)
// df não possui lrenda nem renda2
// df2 as possui
\end{lstlisting}

\section{\texttt{filter}}

Seleciona linhas que satisfazem uma condição:

\begin{lstlisting}
let recentes = filter(df, ano >= 2010)
let ricos = filter(df, renda > 5000 and uf == "SP")
\end{lstlisting}

A forma imperativa usa \lstinline{keep}:

\begin{lstlisting}
keep df if ano >= 2010
\end{lstlisting}

\section{\texttt{select} e \texttt{drop}}

\begin{lstlisting}
// mantém só essas colunas
let slim = select(df, ano, uf, renda)
let sem_id = drop(df, id, codigo)        // remove essas colunas
\end{lstlisting}

\section{\texttt{rename}}

\begin{lstlisting}
rename df antiga nova
// renomeia renda_familiar → renda
rename df renda_familiar renda
\end{lstlisting}

\section{\texttt{sort}}

\begin{lstlisting}
sort df ano               // ascendente
sort df renda desc        // descendente
sort df uf ano desc       // múltiplas colunas
\end{lstlisting}

\section{\texttt{dropna}}

Remove linhas com valores ausentes:

\begin{lstlisting}
dropna df                 // remove linhas com qualquer NaN
// remove linhas com NaN em renda ou gini
dropna df renda gini
\end{lstlisting}

\section{\texttt{append}}

Empilha dois DataFrames verticalmente (equivalente a \texttt{rbind}):

\begin{lstlisting}
let df_total = append(df_2024, df_2025)
\end{lstlisting}

As colunas devem ser compatíveis. Colunas ausentes em um dos DataFrames são
preenchidas com \texttt{NaN}.

\section{\texttt{collapse} e \texttt{group\_by}}

\lstinline{group_by} agrega por grupos:

\begin{lstlisting}[caption={Agregação por grupo}]
let media_uf = group_by(df, uf,
  renda_media = mean(renda),
  n           = count()
)
\end{lstlisting}

\lstinline{collapse} é a forma imperativa estilo Stata:

\begin{lstlisting}
collapse df (mean) renda gini, by(uf ano)
\end{lstlisting}

\section{\texttt{pivot\_longer} e \texttt{pivot\_wider}}

\begin{lstlisting}[caption={Formato largo para longo}]
let longo = pivot_longer(df,
  cols  = ["t1", "t2", "t3"],
  names = "trimestre",
  values = "valor"
)
\end{lstlisting}

\begin{lstlisting}[caption={Formato longo para largo}]
let largo = pivot_wider(df,
  id     = "municipio",
  names  = "ano",
  values = "pib"
)
\end{lstlisting}

\section{\texttt{replace}}

Substitui valores condicionalmente:

\begin{lstlisting}
// zera valores negativos
replace df renda = 0 if renda < 0
replace df gini = . if gini > 1         // marca como missing
\end{lstlisting}

\section{\texttt{tabulate}}

Tabela de frequências:

\begin{lstlisting}
tabulate df uf
tabulate df uf ano    // tabulação cruzada
\end{lstlisting}
